\ifdefined\gkver\else\def\gkver{student}\fi
\documentclass[\gkver]{gaokaozhenti}
\begin{document}

\chapter{2011年重庆理综}
%% paperType: 理综物理部分

\begin{enumerate}
\item
%% number: 14
%% paperName: 2011年重庆理综第 14 题 6 分
%% typeId: 1
%% type: 单选题
%% chapter: 直线运动
%% point: 自由落体运动
%% method: 无
%% score: 6
%% degree: 650
%% duplicateId: 0
%% body:
某人估测一竖直枯井深度，从井口静止释放一石头并开始计时，经$2\Us$听到石头落地声，由此可知井深约为（不计声音传播时间，重力加速度$g$取$10\Umsq$）\xzanswer{B}
\fourchoices[answer=B]
{$10\Um$}
{$20\Um$}
{$30\Um$}
{$40\Um$}
%% answer: B
%% memo:
\memoanswer{
石头做自由落体运动，根据位移公式 $h=\frac{1}{2}gt^{2}=0.5\times10\times4\Um=20\Um$。答案 B。
}

\item
%% number: 15
%% paperName: 2011年重庆理综第 15 题 6 分
%% typeId: 1
%% type: 单选题
%% chapter: 气体性质
%% point: 热力学第一定律
%% method: 无
%% score: 6
%% degree: 650
%% duplicateId: 0
%% body:
某汽车后备箱内安装有撑起箱盖的装置，它主要由汽缸和活塞组成。开箱时，密闭于汽缸内的压缩气体膨胀，将箱盖顶起，如图所示。在此过程中，若缸内气体与外界无热交换，忽略气体分子间相互作用，则缸内气体\xzanswer{A}
\onepicture[w=5cm]{\includesvg{figs/15}}
\fourchoices[answer=A]
{对外做正功，分子的平均动能减小}
{对外做正功，内能增大}
{对外做负功，分子的平均动能增大}
{对外做负功，内能减小}
%% answer: A
%% memo:
\memoanswer{
密闭于气缸内的压缩气体膨胀对外做正功 $W<0$，缸内气体与外界无热交换说明 $Q=0$，忽略气体分子间相互作用，说明内能是所有分子动能的总和。根据热力学第一定律 $\Delta U=W+Q$，可知内能减小，分子平均动能减小，温度降低。所以只有 A 正确。
}

\item
%% number: 16
%% paperName: 2011年重庆理综第 16 题 6 分
%% typeId: 1
%% type: 单选题
%% chapter: 原子和原子核
%% point: 天然放射性
%% method: 无
%% score: 6
%% degree: 650
%% duplicateId: 0
%% body:
核电站核泄漏的污染物中含有碘131和铯137。碘131的半衰期约为8天，会释放$\beta$射线；铯137是铯133的同位素，半衰期约为30年，发生衰变期时会辐射$\gamma$射线。下列说法正确的是\xzanswer{D}
\fourchoices[answer=D]
{碘131释放的$\beta$射线由氦核组成}
{铯137衰变时辐射出的$\gamma$光子能量小于可见光光子能量}
{与铯137相比，碘131衰变更慢}
{铯133和铯137含有相同的质子数}
%% answer: D
%% memo:
\memoanswer{
$\beta$ 射线是高速电子流，$\alpha$ 射线才是由氦核组成，A 错误；$\gamma$ 光子在所有电磁波中频率最高，能量最大，B 错误；半衰期越小衰变越快，应该是碘 131 衰变更快，C 错误；铯 133 和铯 137 是同位素，质子数相同中子数不同，D 正确。
}

\item
%% number: 17
%% paperName: 2011年重庆理综第 17 题 6 分
%% typeId: 1
%% type: 单选题
%% chapter: 机械波
%% point: 波与振动图象
%% method: 无
%% score: 6
%% degree: 450
%% duplicateId: 0
%% body:
介质中坐标原点$O$处的波源在$t=0$时刻开始振动，产生的简谐波沿$x$轴正向传播，$t_{0}$时刻传到$L$处，波形如图所示。下列能描述$x_{0}$处质点振动的图象是\xzanswer{C}
\onepicture[w=5.5cm]{figs/17.png}
\fourchoices[answer=C, ispicture=true, h=3cm]
{figs/17a.png}
{figs/17b.png}
{figs/17c.png}
{figs/17d.png}
%% answer: C
%% memo:
\memoanswer{
从波形图可以看出，$t_0$ 时刻传到 $L=3\lambda$ 处，说明 $t_0=3T$。简谐波沿 $x$ 轴正向传播，则在 $t_0$ 时刻 $x_0$ 处质点的运动方向和波源在 $t=0$ 时刻的运动方向相同，是沿 $y$ 轴的负方向的，即每一个质点的起振方向都是沿 $y$ 轴的负方向的，则 C、D 可能正确。由于 $\lambda<x_0<5\lambda/4$，说明在 $T<t<5T/4$ 的时间内 $x_0$ 处质点没有振动，所以在 $t_0$ 时刻 $x_0$ 处质点的振动时间是 $3T-5T/4<t_0-t<3T-T$，即 $7T/4<t_0-t<2T$，即振动图像中 $t_0$ 时刻前有少于 2 个多于 7/4 个的完整图形，所以 C 正确。
}

\item
%% number: 18
%% paperName: 2011年重庆理综第 18 题 6 分
%% typeId: 1
%% type: 单选题
%% chapter: 光学
%% point: 光的折射
%% method: 无
%% score: 6
%% degree: 450
%% duplicateId: 0
%% body:
在一次讨论中，老师问道：“假如水中相同深度处有 a、b、c 三种不同颜色的单色点光源，有人在水面上方同等条件下观测发现，b 在水下的像最深，c 照亮水面的面积比 a 的大。关于这三种光在水中的性质，同学们能做出什么判断？”有同学回答如下：①c 光的频率最大；②a 光的传播速度最小；③b 光的折射率最大；④a 光的波长比 b 光的短。根据老师的假定，以上回答正确的是\xzanswer{C}
\fourchoices[answer=C]
{①②}
{①③}
{②④}
{③④}
%% answer: C
%% memo:
\memoanswer{
根据视深公式 $h'=\frac{h}{n}$ 说明频率最小的光，水对它的折射率最小，在水下的像最深，所以 b 的折射率最小，频率最小，波长最大，传播速度最大，③错误，④正确；照亮水面的圆面积的半径 $R$ 与临界角 $C$ 满足 $\tan C=\frac{R}{h}$，又 $\sin C=\frac{1}{n}$，c 照亮水面的面积比 a 的大，则 c 的临界角大，水对 c 的折射率小，所以 a 的折射率最大，a 的频率最大，a 的传播速度最小，①错误，②正确。所以正确的说法是②和④两项，选择 C。
}

\item
%% number: 19
%% paperName: 2011年重庆理综第 19 题 6 分
%% typeId: 1
%% type: 单选题
%% chapter: 电场
%% point: 电场叠加
%% method: 无
%% score: 6
%% degree: 450
%% duplicateId: 0
%% body:
如图所示，电量为$+q$和$-q$的点电荷分别位于正方体的顶点，正方体范围内电场强度为零的点有\xzanswer{D}
\onepicture[w=4cm]{figs/19.png}
\fourchoices[answer=D]
{体中心、各面中心和各边中点}
{体中心和各边中点}
{各面中心和各边中点}
{体中心和各面中心}
%% answer: D
%% memo:
\memoanswer{
两个等量同种电荷在其连线的中点处的合场强为零。两个等量同种正电荷在其连线的中垂线上的合场强沿中垂线指向远离正电荷的方向。两个等量同种负电荷在其连线的中垂线上的合场强沿中垂线指向负电荷的方向。在正方体的上面中心，上面的四个电荷分成两组产生的场强都是零，下面的四个电荷分成两组产生的场强等大反向，所以正方体的上面中心处的合场强为零，同理所有各面中心处的合场强都为零。在体中心，可以将八个电荷分成四组，产生的合场强为零。而在各边中心，场强无法抵消，合场强不为零。正确答案是 D。
}

\item
%% number: 20
%% paperName: 2011年重庆理综第 20 题 6 分
%% typeId: 1
%% type: 单选题
%% chapter: 电路
%% point: 伏安法测电阻
%% method: 无
%% score: 6
%% degree: 450
%% duplicateId: 0
%% body:
在测量电珠伏安特性实验中，同学们连接的电路中有四个错误电路，如图所示。电源内阻不计，导线连接良好。若将滑动变阻器的触头置于左端，闭合 S，在向右端滑动触头过程中，会分别出现如下四种现象：（A）电珠$L$不亮；电流表示数几乎为零；（B）电珠$L$亮度增加；电流表示数增大；（C）电珠$L$开始不亮；后来忽然发光；电流表从示数不为零到线圈烧断；（D）电珠$L$不亮；电流表从示数增大到线圈烧断。与上述 abcd 四种现象对应的电路序号为\xzanswer{A}
\onepicture[w=13cm]{figs/20.png}
\fourchoices[answer=A]
{③①②④}
{③④②①}
{③①④②}
{②①④③}
%% answer: A
%% memo:
\memoanswer{
在①中伏特表和安培表的连接是错误的，将滑动变阻器的触头置于左端，闭合 S 后电珠 $L$ 不亮，在向右端滑动触头过程中，电珠 $L$ 逐渐变亮，由于总电阻减小，干路电流增大，电流表示数增大，是 b；在②中安培表连接错误，应该连接在干路上，开始时滑动变阻器的电阻最大，通过电珠的电流太小，电珠 $L$ 不亮，在向右端滑动触头过程中，电珠后来忽然发光，由于大部分电流流过电流表造成电流表从示数不为零到线圈烧断，是 c；在③中伏特表和安培表应该互换位置，由于伏特表串联造成电珠 $L$ 不亮，电流表示数几乎为零，是 a；在④中安培表使电珠短路，电珠 $L$ 始终不亮，电流表从示数增大到线圈烧断，是 d。即①②③④分别对应 bcad。正确的是 A。
}

\item
%% number: 21
%% paperName: 2011年重庆理综第 21 题 6 分
%% typeId: 1
%% type: 单选题
%% chapter: 曲线运动
%% point: 开普勒定律
%% method: 无
%% score: 6
%% degree: 450
%% duplicateId: 0
%% body:
某行星和地球绕太阳公转的轨道均可视为圆。每过 $N$ 年，该行星会运行到日地连线的延长线上，如图所示。该行星与地球的公转半径比为\xzanswer{B}
\onepicture[w=4cm]{figs/21.png}
\fourchoices[answer=B]
{$\left(\frac{N+1}{N}\right)^{\frac{2}{3}}$}
{$\left(\frac{N}{N-1}\right)^{\frac{2}{3}}$}
{$\left(\frac{N+1}{N}\right)^{\frac{3}{2}}$}
{$\left(\frac{N}{N-1}\right)^{\frac{3}{2}}$}
%% answer: B
%% memo:
\memoanswer{
由图可知行星的轨道半径大周期长。每过 $N$ 年，该行星会运行到日地连线的延长线上，说明从最初在日地连线的延长线上开始，每一年地球都在行星的前面比行星多转圆周的 $N$ 分之一，$N$ 年后地球转了 $N$ 圈，比行星多转 1 圈，即行星转了 $N-1$ 圈从而再次在日地连线的延长线上。所以行星的周期是 $\frac{N}{N-1}$ 年，根据开普勒第三定律有 $\frac{r_\text{地}^{3}}{r_\text{行}^{3}}=\frac{T_\text{地}^{2}}{T_\text{行}^{2}}$，则答案是 B。
}

\item
%% number: 22
%% paperName: 2011年重庆理综第 22 题 13 分
%% typeId: 4
%% type: 实验题
%% chapter: 运动和力
%% point: 验证牛顿第二定律
%% method: 无
%% score: 13
%% degree: 450
%% duplicateId: 0
%% body:
\textbf{（1）}（6 分）某电动机的三组线圈①、②、③阻值相同，均为几欧姆，接法可能是图甲、图乙两种之一，A、B 和 C 是外接头。现有一组线圈断路，维修人员通过多用电表测量外接头之间的电阻来判断故障，若测量 A 和 B 之间、B 和 C 之间、A 和 C 之间的电阻时，多用电表指针偏转分别如图\ref{2011重庆理综22c}、\ref{2011重庆理综22d}、\ref{2011重庆理综22e}所示，则测量中使用的欧姆档的倍率是\tkanswer{$\times$1}（填 $\times$1、$\times$10、$\times$100 或 $\times$1k），三组线圈的接法是\tkanswer{乙}（填甲或乙），断路线圈是\tkanswer{③}（填①、②或③）。

\twopicture[numA=甲, numB=乙, labelA=2011重庆理综22a, labelB=2011重庆理综22b, align=right, hA=3cm, hB=3cm]
{figs/22a.png}{figs/22b.png}
\threepicture[numA={（a）}, numB={（b）}, numC={（c）}, labelA=2011重庆理综22c, labelB=2011重庆理综22d, labelC=2011重庆理综22e, align=right, hA=2cm, hB=2cm, hC=2cm]
{figs/22c.png}{figs/22d.png}{figs/22e.png}

\textbf{（2）}（13 分）某同学设计了图\ref{2011重庆理综22f}所示的装置，利用米尺、秒表、轻绳、轻滑轮、轨道、滑块、托盘和砝码等器材来测定滑块和轨道间的动摩擦因数$\mu$。滑块和托盘上分别放有若干砝码，滑块质量为$M$，滑块上砝码总质量为$m'$，托盘和盘中砝码的总质量为$m$。实验中，滑块在水平轨道上从 A 到 B 做初速为零的匀加速直线运动，重力加速度$g$取$10\Umsq$。
\begin{enumerate}
	\item
	为测量滑块的加速度$a$，须测出它在 A、B 间运动的\tkanswer{位移}与\tkanswer{时间}，计算$a$的运动学公式是\tkanswer{$a=\frac{2s}{t^{2}}$}；
	\item
	根据牛顿运动定律得到$a$与$m$的关系为：$a=\frac{(1+\mu)g}{M+(m+m')}m-\mu g$，他想通过多次改变$m$，测出相应的$a$值，并利用上式来计算$\mu$。若要求$a$是$m$的一次函数，必须使上式中的\tkanswer{$m+m'$}保持不变，实验中应将从托盘中取出的砝码置于\tkanswer{滑块上}；
	\item
	实验得到$a$与$m$的关系如图\ref{2011重庆理综22g}所示，由此可知$\mu=$\tkanswer{0.23}（取两位有效数字）。
\end{enumerate}
\onepicture[num=1, label=2011重庆理综22f, align=right, w=7cm]{figs/22f.png}
\onepicture[num=2, label=2011重庆理综22g, align=right, w=6cm]{figs/22g.png}
%% answer:
\jdanswer{
\begin{enumerate}
	\item
	位移，时间，$a=\frac{2s}{t^{2}}$

	\item
	$m+m'$，滑块上

	\item
	0.23（0.21$\sim$0.25）

\end{enumerate}
}
%% memo:
\memoanswer{
（1）由于所测电阻均为几欧姆，串联后还是几欧姆，所以欧姆档的倍率选择 $\times$1；所测值分别是 $4\UO$、$4\UO$ 和 $8\UO$，说明接法是乙，断路的是③。因为采用甲接法时测量值应该有两个是无穷大，一个是 $8\UO$。

（2）①滑块在水平轨道上从 A 到 B 做初速为零的匀加速直线运动，根据 $s=\frac{1}{2}at^{2}$ 得 $a=\frac{2s}{t^{2}}$，所以需要测量的是位移和时间。

②根据整体法有

$a=\frac{mg-\mu(M+m')g}{M+m+m'}=\frac{mg-\mu(M+m+m')g+\mu mg}{M+m+m'}=\frac{(1+\mu)g}{M+m+m'}m-\mu g$

若要求 $a$ 是 $m$ 的一次函数必须使 $\frac{(1+\mu)g}{M+m+m'}$ 不变，即使 $m+m'$ 不变，在增大 $m$ 时等量减小 $m'$，所以实验中应将从托盘中取出的砝码置于滑块上。

③将 $\frac{(1+\mu)g}{M+m+m'}$ 取为 $k$，有 $k=\frac{a_{1}+\mu g}{m_{1}}=\frac{a_{2}+\mu g}{m_{2}}$，在图像上取两点将坐标代入解得 $\mu=0.23$（在 0.21 到 0.25 之间是正确的）。
}

\item
%% number: 23
%% paperName: 2011年重庆理综第 23 题 16 分
%% typeId: 6
%% type: 计算题
%% chapter: 电磁感应
%% point: 电磁感应能量分析
%% method: 无
%% score: 16
%% degree: 650
%% duplicateId: 0
%% body:
有人设计了一种可测速的跑步机，测速原理如图所示，该机底面固定有间距为$L$、长度为$d$的平行金属电极。电极间充满磁感应强度为$B$、方向垂直纸面向里的匀强磁场，且接有电压表和电阻$R$。绝缘橡胶带上镀有间距为$d$的平行细金属条，磁场中始终仅有一根金属条，且与电极接触良好，不计金属电阻，若橡胶带匀速运动时，电压表读数为$U$，求：
\begin{enumerate}
	\item
	橡胶带匀速运动的速率；
	\item
	电阻$R$消耗的电功率；
	\item
	一根金属条每次经过磁场区域克服安培力做的功。
\end{enumerate}
\onepicture[align=right, w=7cm]{figs/23.png}
%% answer:
\jdanswer{
\begin{enumerate}
	\item
	$v=\frac{U}{BL}$

	\item
	$P=\frac{U^{2}}{R}$

	\item
	$W=\frac{BLUd}{R}$

\end{enumerate}
}
%% memo:
\memoanswer{
（1）设电动势为 $E$，橡胶带运动速率为 $v$。

由：$E=BLv$，$E=U$

得：$v=\frac{U}{BL}$

（2）设电功率为 $P$，则 $P=\frac{U^{2}}{R}$

（3）设电流强度为 $I$，安培力为 $F$，克服安培力做的功为 $W$。

$I=\frac{U}{R}$，$F=BIL$，$W=Fd$

得：$W=\frac{BLUd}{R}$
}

\item
%% number: 24
%% paperName: 2011年重庆理综第 24 题 18 分
%% typeId: 6
%% type: 计算题
%% chapter: 运动和力
%% point: 动量和能量
%% method: 无
%% score: 18
%% degree: 450
%% duplicateId: 0
%% body:
如图所示，静置于水平地面的三辆手推车沿一直线排列，质量均为$m$，人在极短时间内给第一辆车一水平冲量使其运动，当车运动了距离$L$时与第二辆车相碰，两车以共同速度继续运动了距离$L$时与第三车相碰，三车以共同速度又运动了距离$L$时停止。车运动时受到的摩擦阻力恒为车所受重力的$k$倍，重力加速度为$g$，若车与车之间仅在碰撞时发生相互作用，碰撞时间很短，忽略空气阻力，求：
\begin{enumerate}
	\item
	整个过程中摩擦阻力所做的总功；
	\item
	人给第一辆车水平冲量的大小；
	\item
	第一次与第二次碰撞系统功能损失之比。
\end{enumerate}
\onepicture[align=right, w=6.5cm]{figs/24.png}
%% answer:
\jdanswer{
\begin{enumerate}
	\item
	$W=-6kmgL$

	\item
	$I=2m\sqrt{7kgL}$

	\item
	$\Delta E_{k1}/\Delta E_{k2}=13/3$

\end{enumerate}
}
%% memo:
\memoanswer{
（1）设运动过程中摩擦阻力做的总功为 $W$，则

$W=-kmgL-2kmgL-3kmgL=-6kmgL$

（2）设第一车初速度为 $u_{0}$，第一次碰前速度为 $v_{1}$，碰后共同速度为 $u_{1}$；第二次碰前速度为 $v_{2}$，碰后共同速度为 $u_{2}$；人给第一车的水平冲量大小为 $I$。

由：$-kmgL=\frac{1}{2}mv_{1}^{2}-\frac{1}{2}mu_{0}^{2}$

$-2kmgL=\frac{1}{2}(2m)v_{2}^{2}-\frac{1}{2}(2m)u_{1}^{2}$

$-3kmgL=0-\frac{1}{2}(3m)u_{2}^{2}$

$mv_{1}=2mu_{1}$

$2mv_{2}=3mu_{2}$

得：$I=mu_{0}-0=2m\sqrt{7kgL}$

（3）设两次碰撞中系统动能损失分别为 $\Delta E_{k1}$ 和 $\Delta E_{k2}$

由 $\Delta E_{k1}=\frac{13}{2}kmgL$，$\Delta E_{k2}=\frac{3}{2}kmgL$

得：$\Delta E_{k1}/\Delta E_{k2}=13/3$
}

\item
%% number: 25
%% paperName: 2011年重庆理综第 25 题 19 分
%% typeId: 6
%% type: 计算题
%% chapter: 磁场
%% point: 带电粒子在磁场中的运动
%% method: 无
%% score: 19
%% degree: 250
%% duplicateId: 0
%% body:
某仪器用电场和磁场来控制电子在材料表面上方的运动。如图所示，材料表面上方矩形区域PP$'$N$'$N充满竖直向下的匀强电场，宽为$d$；矩形区域NN$'$M$'$M充满垂直纸面向里的匀强磁场，磁感应强度为$B$，长为3$s$，宽为$s$；NN$'$为磁场与电场之间的薄隔离层。一个电荷量为$e$、质量为$m$、初速为零的电子，从P点开始被电场加速经隔离层垂直进入磁场，电子每次穿越隔离层，运动方向不变，其动能损失是每次穿越前动能的10\%，最后电子仅能从磁场边界M$'$N$'$飞出。不计电子所受重力。
\begin{enumerate}
	\item
	求电子第二次与第一次圆周运动半径之比；
	\item
	求电场强度的取值范围；
	\item
	A是M$'$N$'$的中点，若要使电子在A、M$'$间垂直于AM$'$飞出，求电子在磁场区域中运动的时间。
\end{enumerate}
\onepicture[align=right, w=7cm]{figs/25.png}
%% answer:
\jdanswer{
\begin{enumerate}
	\item
	$R_{1}:R_{2}=0.9$

	\item
	$\frac{B^{2}es^{2}}{80md}<E\leq\frac{5B^{2}es^{2}}{9md}$

	\item
	$t=\frac{5\pi m}{2eB}$

\end{enumerate}
}
%% memo:
\memoanswer{
（1）设圆周运动的半径分别为 $R_1$、$R_2$、$\ldots$、$R_n$、$R_{n+1}$，第一和第二次圆周运动速率分别为 $v_1$ 和 $v_2$，动能分别为 $E_{k1}$ 和 $E_{k2}$

由：$E_{k2}=0.81E_{k1}$，$R_{1}=\frac{mv_{1}}{Be}$，$R_{2}=\frac{mv_{2}}{Be}$，$E_{k1}=\frac{1}{2}mv_{1}^{2}$，$E_{k2}=\frac{1}{2}mv_{2}^{2}$，得：$R_{1}:R_{2}=0.9$

（2）设电场强度为 $E$，第一次到达隔离层前的速率为 $v'$

由：$eEd=\frac{1}{2}mv'^2$，$0.9\times\frac{1}{2}mv'^2=\frac{1}{2}mv_1^2$，$R_1\leq s$

得：$E\leq\frac{5B^{2}es^{2}}{9md}$

又由：$R_{n}=0.9^{n-1}R_{1}$，$2R_{1}(1+0.9+0.9^{2}+\ldots+0.9^{n}+\ldots)=3s$

得：$E>\frac{B^{2}es^{2}}{80md}$

电场强度的取值范围为 $\frac{B^{2}es^{2}}{80md}<E\leq\frac{5B^{2}es^{2}}{9md}$

（3）设电子在匀强磁场中圆周运动的周期为 $T$，运动的半圆周个数为 $n$，运动总时间为 $t$

由题意，有：$\frac{2R_{1}(1-0.9^{n})}{1-0.9}+R_{n+1}=3s$，$R_{1}\leq s$，$R_{n+1}=0.9^{n}R_{1}$，$R_{n+1}\geq\frac{s}{2}$

得：$n=2$

由：$T=\frac{2\pi m}{eB}$

得：$t=\frac{5\pi m}{2eB}$

轨迹如图所示，电子在磁场区域中运动的时间为 $1.25T$。

\onepicture[w=8cm]{figs/25_an.png}
}

\end{enumerate}

\end{document}
